\section{A constant approximation for expected depth}
\label{sec:mass-improvement}

The expected-depth algorithm distinguishes balanced and flat states
by probability mass. Like the maximum-depth algorithm, it learns
the complete hull pattern before recursing. Its discovery ordering
uses exact weighted coverage rather than the heights of child codes.
The proof uses the same three steps in this section and in all later
expected-depth improvements: compare discovery with queries deleted
from an optimal witness path; bound the information still needed to
recover the full pattern; then telescope these two local credits.
Lemma~\ref{lem:stage-conversion} performs the last step once for all
variants. The algorithm itself uses only coverage, endpoint codes, and
rational scores; the witness path and information quantities are proof
objects.
Fix
\[
 \beta=\frac{3512863}{10^6},\qquad
 \lambda=\frac{\beta+1}{2\ln\beta},\qquad
 \nu=1+\frac1{\ln\beta},
\]
and the accounting coefficients
\[
 A_\beta=2\lambda,\qquad
 C_\beta=\frac{\Phi_\beta(1)+\Phi_\beta(2)}2,\qquad
 B_\beta=4+C_\beta.
\]
For this $\beta$, $\lfloor\beta\rfloor=3$ and
$\lfloor2\beta\rfloor=7$. Splitting the integrals at their integer
breakpoints gives
\[
 C_\beta=3-\frac{\ln(945/8)}{4\ln\beta},\qquad
 B_\beta=7-\frac{\ln(945/8)}{4\ln\beta}.
\]
The chosen $\beta$ satisfies $\lambda>\nu$.

Revealing a pattern makes an initial segment of an optimal witness
path unnecessary; those deleted queries pay for its discovery.
Entropy pays for the pattern information that remains. We seek a
local bound of the form
\[
 \text{expected queries in a stage}
 \ \le\ A_\beta\Delta_{\mathrm E}(S,\Pi)+cI(S,\Pi)
\]
for a constant $c$ and the stage's output partition $\Pi$.
Unlike a bound proportional to $\OPT_{\mathrm E}(S)$ alone, this inequality sums
over nested states without charging the same optimum repeatedly.

\subsection{Hull tests and their answer patterns}

Fix a strict test witness $p$ and use the common notation of
Definitions~\ref{def:normalized-queries} and~\ref{def:hull-pattern}.
These hull vertices are used for discovery, not as a replacement for
the original query family. Once the pattern class is known, the
recursive algorithm again has access to every supplied line.

Lemma~\ref{lem:hull-structure} gives cap coverage and interval
patterns. Lemma~\ref{lem:mssc} gives a deterministic ordering with
latency and length bounds. Endpoint recovery is described in
Section~\ref{sec:endpoint-geometry}.

\subsection{The queries a pattern makes unnecessary}

\begin{lemma}[Supporting-prefix credit]\label{lem:support-prefix}
Let $T^*$ be an expected-depth-optimal proper tree for $S$.
Follow $p$ through it, with queries $Q=(i_1,\ldots,i_q)$ and terminal
hypothesis $z$.
For $A\in\Pi$, let $t_A$ be the first index $j$ for which the pattern
of $A$ is positive at a vertex of $J_{i_j}$; set $t_A=q+1$ if there
is no such index.
Then
\begin{equation}\label{eq:support-prefix-credit}
 \Delta_{\mathrm E}(S,\Pi)\ge
 \sum_{A\in\Pi}w(A)(t_A-1).
\end{equation}
Every class with $t_A=q+1$ is the singleton $\{z\}$.
If $A_0$ exists, then $z$ is its member and $t_{A_0}=q+1$.
\end{lemma}
\begin{proof}
For $j<t_A$, every supporting hull bit in $J_{i_j}$ is negative on
the whole class $A$. Equation~\eqref{eq:hull-cap-cover} forces
test $i_j$ to have the witness's answer throughout $A$.
Inductively, all members of $A$ follow the witness path through
its first $t_A-1$ queries. Restricting $T^*$ to $A$ contracts
those queries and reduces its weighted depth by
$w(A)(t_A-1)$. Other contractions can only reduce it further.

If there is no supporting hit, all members of $A$ follow the
entire path to leaf $z$, so $A=\{z\}$. The all-zero class always
does this. Summing the reductions over disjoint classes gives
\[
 \Delta_{\mathrm E}(S,\Pi)\ge
 c_{\mathrm E}(T^*)-\sum_{A\in\Pi}c_{\mathrm E}(T^*|_A),
\]
and the right side is at least the stated prefix sum.
The inequality uses $\OPT_{\mathrm E}(A)\le c_{\mathrm E}(T^*|_A)$: each restricted tree
is a valid identifying tree for its class.
\end{proof}

The first supporting hit need not be the original test's first
positive answer. That is harmless: the proof only needs the
stronger implication that all-negative supports force a constant
original answer. The algorithm will query the supports directly.

\subsection{Finding a positive hull answer}

Compute the pattern classes and put $U=S\setminus A_0$.
Hull caps cover $U$. Apply Lemma~\ref{lem:mssc} to its weighted
interval demands with $p_0=w(A_0)$. Query the resulting ordering
until a positive answer or a singleton can be inferred.
If every answer is negative and
$A_0=\{z_0\}$, completing the cover identifies $z_0$.
When $A_0$ is absent, a complete all-negative route is impossible.

The shift is chosen by the joint first-cover and waiting-time
objective, using the actual all-zero weight rather than a guessed
terminal point. Its average waiting-time bound is sharper than
the uniform length bound. Choosing the joint objective makes that
saving available to the same deterministic schedule.

To analyze discovery, follow a witness through an optimal expected-
depth tree and replace each query by its one or two hull supports.
This is only a comparison ordering; the algorithm cannot see it.
Within a support block, querying the heavier newly hit part first
ensures that at least half of its mass is found at the earlier of
the two possible positions. The exact batch analysis can then charge
the actual schedule to those integer first-cover times.

\begin{lemma}[Discovery bound]\label{lem:hull-discovery}
The weighted expected number $W_{\mathrm{disc}}(S)$ of discovery
queries satisfies
\[
 W_{\mathrm{disc}}(S)\le
 A_\beta\Delta_{\mathrm E}(S,\Pi)+(B_\beta-4)b(S).
\]
\end{lemma}
\begin{proof}
Singleton states require no query. Otherwise $U$ is nonempty,
because the all-zero class has at most one hypothesis.
Use the witness path and supporting-prefix indices of
Lemma~\ref{lem:support-prefix}. Put
\[
 D_+=\sum_{A\in\Pi,\ A\ne A_0}w(A)(t_A-1).
\]
Expand each path query into its at most two real supports.
Within that block, put first the support hitting the larger
weight of classes whose first supporting index is this block.
It hits at least half of that newly supported mass.
Their preceding blocks use at most $2(t_A-1)$ tests.
Consequently at least half of a block's newly supported mass is
hit by position $2t_A-1$, and all of it by position $2t_A$.
Removing repetitions only improves these positions.
This mass ordering belongs to the comparison proof, not to the
algorithm, which does not know the optimal witness path.

The discrete curve must be applied before replacing these two
positions by their average. Write $\eta=\eta_\beta$ and put
\[
 g(t)=\lambda t+\eta-\Phi_\beta(t)
     =\frac1{2\ln\beta}\int_1^t
       \frac{(x-\lfloor x\rfloor)
              +(\beta x-\lfloor\beta x\rfloor)}x\,dx.
\]
The function $g$ is the exact saving lost by replacing the discrete
curve $\Phi_\beta$ with its coarser linear upper bound
$\lambda t+\eta$. Its integral representation shows that it is
nonnegative and nondecreasing, with $g(1)=0$.
Define
\[
 \varepsilon_\beta(t)
 =\frac{g(2t-1)+g(2t)-g(2)}2.
\]
It too is nonnegative and nondecreasing, because $g$ is, and
$\varepsilon_\beta(1)=0$.
For every integer $t\ge1$,
\[
 \begin{aligned}
 \frac{\Phi_\beta(2t-1)+\Phi_\beta(2t)}2
 &=2\lambda(t-1)+\frac32\lambda+\eta
     -\frac{g(2t-1)+g(2t)}2\\
 &=2\lambda(t-1)+C_\beta-\varepsilon_\beta(t),
 \end{aligned}
\]
because the last average is at least $g(2)/2$ and
$C_\beta=\frac32\lambda+\eta-g(2)/2$.
Since $\Phi_\beta$ is increasing, the worst placement consistent
with the block order puts half the new mass at position $2t-1$ and
half at position $2t$. Its contribution is therefore bounded by
the displayed average times the block mass. A singleton support
places all of its mass at the earlier position and obeys the same
bound.

If $A_0$ is absent, the only possibly unhit class after all blocks
is the singleton $\{z\}$. Append one positive hull test for it
if needed. Its position is at most $2q+1=2(t_{\{z\}}-1)+1$.
It too obeys the preceding bound.
Let $m_t$ be the mass of classes whose first supporting block is
$t$, including the appended terminal singleton when present.
For the resulting comparison order $Q_{\mathrm{sup}}$,
\[
 \sum_{h\in U}\pi_h\Phi_\beta(t_{Q_{\mathrm{sup}}}(h))
 \le2\lambda D_++C_\beta b(S)
       -\sum_t m_t\varepsilon_\beta(t).
\]
If $A_0=\{z_0\}$, every nonzero class has $t_A\le q$.
The support expansion covers $U$ with at most $2q$ tests.
Its minimum hull cover size $k_0$ is therefore at most $2q$.
Supporting-prefix credit gives
$D_++\pi_{z_0}q\le\Delta_{\mathrm E}$.
When $A_0$ is absent, it instead gives $D_+\le\Delta_{\mathrm E}$.
Apply the joint inequality in Lemma~\ref{lem:mssc} to this
comparison order and the actual zero weight:
\[
\begin{aligned}
 W_{\mathrm{disc}}
 &\le2\lambda D_++C_\beta b(S)+w(A_0)Z_\beta(k_0)\\
 &\le2\lambda D_++C_\beta b(S)+2\nu w(A_0)q\\
 &\le A_\beta\Delta_{\mathrm E}(S,\Pi)+(B_\beta-4)b(S).
\end{aligned}
\]
The middle inequality drops the nonpositive correction
$-w(A_0)/\ln\beta$. For the last inequality, if $A_0$ is present,
the support-prefix credit gives
$D_++w(A_0)q\le\Delta_{\mathrm E}$ and $\nu\le\lambda$ gives
$2\lambda D_++2\nu w(A_0)q\le2\lambda\Delta_{\mathrm E}$.
If $A_0$ is absent, $D_+\le\Delta_{\mathrm E}$ directly. Finally,
$A_\beta=2\lambda$ and $B_\beta-4=C_\beta$.
Before discarding these savings, the same calculation gives
\begin{equation}\label{eq:discrete-discovery-credit}
 W_{\mathrm{disc}}\le
 A_\beta\Delta_{\mathrm E}+C_\beta b(S)
 -\sum_t m_t\varepsilon_\beta(t)-\frac{w(A_0)}{\ln\beta}.
\end{equation}
Indeed, $Z_\beta(k_0)\le\nu k_0-1/\ln\beta$ and
$k_0\le2q$ retain the stated zero-target credit.
When $A_0$ is absent, its waiting-time terms vanish, so no
bound on $k_0$ is needed there.
Early singleton inference can only reduce the actual work.
\end{proof}

\subsection{The full-pattern stage}

\begin{lemma}[Full-pattern bound]\label{lem:pruning-macro}
Let $\mathcal R$ be the partition at the end of discovery, including
early singleton inferences, and put $J(S)=I(S,\mathcal R)$.
Let $b_{\mathrm{pin}}$ be the mass of positive exits with at least
one cached negative neighbor of the discovered vertex, and
$b_{\mathrm{both}}$ the mass with both neighbors negative.
The discovery and endpoint procedures form a polynomial-time partial
decision tree whose output classes are exactly $\Pi$. Its weighted
expected number of queries satisfies
\begin{equation}\label{eq:pruning-macro}
 \begin{aligned}
 W_{\mathrm{pat}}(S)\le{}&
 A_\beta\Delta_{\mathrm E}(S,\Pi)+B_\beta b(S)
       -2b_{\mathrm{pin}}-2b_{\mathrm{both}}\\
 &+\frac{I(S,\Pi)-J(S)}{\ln2}.
 \end{aligned}
\end{equation}
\end{lemma}
\begin{proof}
Discovery depends only on the full hull pattern. An inferred singleton
therefore has a unique pattern; on other positive exits the endpoints
identify it. The completed negative cover identifies the zero singleton.
Hence the output partition is exactly $\Pi$.

Use Algorithm~\ref{alg:endpoint-code} with the optimal expected-depth
alphabetic codes of Algorithm~\ref{alg:alphabetic-code}. An unpinned exit
uses two searches, a singly pinned exit one, and a doubly pinned exit
none. Their two-query code overheads sum to
$4b-2b_{\mathrm{pin}}-2b_{\mathrm{both}}$; the doubly pinned mass is
also included in $b_{\mathrm{pin}}$.

Let $\Sigma$ denote the full pattern and $Z$ the discovery transcript
under the conditional hypothesis distribution on $S$. The transcript
is determined by $\Sigma$, and after each transcript the realized
endpoint pair $(L,R)$ uniquely describes the positive pattern. Use a
fixed dummy pair on already determined routes. The entropy chain rule
then gives
\[
 w(S)H((L,R)\mid Z)=w(S)H(\Sigma\mid Z)=I(S,\Pi)-J(S),
\]
and $H((L,R)\mid Z)=H(L\mid Z)+H(R\mid L,Z)$.
Lemma~\ref{lem:alphabetic-endpoints} bounds the conditional search cost
by these entropy terms divided by $\ln2$, plus the overheads already
counted. Consequently additional endpoint work is at most
\[
 4b-2b_{\mathrm{pin}}-2b_{\mathrm{both}}
     +\frac{I-J}{\ln2}.
\]
Adding Lemma~\ref{lem:hull-discovery} proves the claim.
\end{proof}

\subsection{The recursive algorithm and its telescoping bound}

Fix a rational $0<\delta<1/3$. We use the near-minimizing threshold
\[
 \delta=\frac{1234989569}{10^{10}}=0.1234989569.
\]
For the chosen $\beta$, this rational approximates the minimizing
balance threshold in \eqref{eq:parameter-ratio}. Its factor is certified
below; algorithmic comparisons require only rational arithmetic.
\begin{symboltable}
$\beta,\ \lambda_\beta,\ \nu_\beta$ & Fixed geometric-batch factor and its positive-target and zero-target coefficients.\\
$A_\beta,\ C_\beta,\ B_\beta$ & Discovery, first-two-block, and endpoint-inclusive coefficients.\\
$\delta,\ h(\delta),\ \sigma(\delta)$ & Mass-balance threshold, binary entropy at that threshold, and capped-discovery entropy per unit mass.\\
$N_\delta,\ s_\delta,\ \chi_\beta(\delta)$ & Number of full early mass blocks, final fractional mass, and guaranteed discrete discovery saving.\\
$\Pi,\ \mathcal R,\ J(S),\ I(S,\Pi)$ & Full-pattern partition, discovery-transcript partition, information learned by discovery, and total pattern information.\\
$b(S),\ b_{\mathrm{pin}},\ b_{\mathrm{both}}$ & Nonzero-pattern mass, mass with a cached negative neighbor, and mass with both neighbors negative.\\
$\Delta_{\mathrm E},\ W_{\mathrm{disc}},\ W_{\mathrm{pat}}$ & Restriction credit, weighted discovery-query cost, and weighted full-stage query cost.\\
$B_{\beta,\delta}$ & Final flat-stage coefficient after the guaranteed discovery saving.\\
\end{symboltable}
Algorithm~\ref{alg:expected-tree}, in its \emph{vertex} variant,
gives the recursion. A mass-balanced supplied line is queried directly.
Otherwise compute the common-witness hull and execute the full-pattern
stage of Lemma~\ref{lem:pruning-macro}; recurse on its classes and stop
at singletons. Its schedule comes from Algorithm~\ref{alg:batch-orders}
with score $\MSSC(G)+w(A_0)|G|$. All decisions use the supplied
incidence table and rational probabilities. The optima and restriction
credit occur only in the analysis.

\begin{algorithm}[H]
\caption{\ExpectedTree: shared recursive expected-depth construction}
\label{alg:expected-tree}
\KwIn{State $S$, original priors; variant \emph{vertex}, \emph{block},
\emph{oriented}, or \emph{adaptive}; its specified rational parameters;
scheduler $\mathcal B$, defaulting to \BatchOrders.}
\KwOut{A supplied-query identification tree.}
\If{$|S|\le1$}{\KwRet its leaf\;}
\If{a mass-balanced supplied line $\ell$ exists}{
  \KwRet $\ell$ with children $\ExpectedTree(S_\ell^-,\text{variant})$
  and $\ExpectedTree(S_\ell^+,\text{variant})$\;
}
Compute $\HullData(S)$; put $p_0=w(A_0)$ and $b=w(S)-p_0$\;
Use real hull vertices as candidates in \emph{vertex} mode;
otherwise use edges of the full hull, retaining its origin marker\;
\If{the hull has at most two real extreme vertices}{
  Build their direct pattern stage and go to child attachment\;
}
\ForEach{candidate schedule $Q$ from $\mathcal B$ on the positive-pattern demands}{
  \eIf{variant is \emph{vertex}}{$F(Q)\gets\MSSC(Q)+p_0|Q|$\;}{
    Use fixed block orientations in \emph{block} mode, and smaller exclusive mass first otherwise\;
    Compute each block's newly hit masses $x,y,z$\;
    Set $D=0$ in \emph{block} mode, $D=\max\{x,y\}$ in \emph{oriented} mode,
    or choose the adaptive mode and $D$ by Algorithm~\ref{alg:block-mode}\;
    $F(Q)\gets2\MSSC_{\rm blk}(Q)+2p_0|Q|-\sum D$\;
  }
}
Choose $Q$ minimizing its rational score and build its discovery stage\;
At each positive exit use expected $\EndpointCode$ to learn the full pattern\;
\textbf{Child attachment:} attach $\ExpectedTree(A,\text{variant})$ to each nonzero pattern class\;
Use a zero-query leaf for the zero singleton; omit that route when absent\;
\KwRet the contracted tree\;
\end{algorithm}

The oriented and mixed-prefix theorems use this same recipe with
different rational balance thresholds; their prefix mixtures are
\emph{analysis certificates}, not new scheduling subroutines.
Algorithms~\ref{alg:first-phase-family} and~\ref{alg:two-phase-family}
replace $\mathcal B$ in adaptive mode and return its best rational-score
candidate; the chosen scheduler and parameters are inherited by children.


\begin{lemma}[Capped-rank credit conversion]\label{lem:capped-rank-credit}
Let $w=b+p_0$, with $p_0\ge0$, and let nonnegative witness masses
$m_t$ satisfy $\sum_t m_t=b$ and $m_t\le2\delta w$.
Suppose coefficients $a_t$ are at least $\tau$ except on a finite
set $E$. For $\mu\ge0$, define
\[
 \chi=\tau-2\delta\sum_{t\in E}(\tau-a_t)_+-\delta\mu,
 \qquad (x)_+=\max\{x,0\}.
\]
If $\chi+\zeta\ge\tau$, then
\begin{equation}\label{eq:capped-rank-conversion}
 \sum_t m_t a_t-\delta w\mu+\zeta p_0\ge\chi b.
\end{equation}
\end{lemma}
\begin{proof}
Charge only the coefficients below $\tau$, each using its individual
mass cap:
\[
 \sum_t m_t a_t\ge\tau b-
     2\delta w\sum_{t\in E}(\tau-a_t)_+.
\]
After subtracting the intercept and adding the zero-route credit,
the right side is
\[
 \tau b-(\tau-\chi)w+\zeta p_0
 =\chi b+p_0(\chi+\zeta-\tau)\ge\chi b.
\]
Thus the same calculation handles arbitrary positive mass $b$ and
arbitrary zero-pattern mass $p_0$; it never assumes $b=w$.
\end{proof}

\begin{lemma}[Flat-state discrete credit]\label{lem:flat-discrete-credit}
Put
\[
 \begin{aligned}
 N_\delta&=\lfloor1/(2\delta)\rfloor,\qquad
 s_\delta=1-2\delta N_\delta,\\
 \chi_\beta(\delta)
 &=2\delta\sum_{t=1}^{N_\delta}\varepsilon_\beta(t)
       +s_\delta\varepsilon_\beta(N_\delta+1),\\
 B_{\beta,\delta}&=4+C_\beta-\chi_\beta(\delta).
 \end{aligned}
\]
If
\begin{equation}\label{eq:flat-discrete-condition}
 \varepsilon_\beta(N_\delta+1)
 \le\chi_\beta(\delta)+\frac1{\ln\beta},
\end{equation}
then a mass-flat stage satisfies
\[
 \begin{aligned}
 W_{\mathrm{pat}}(S)\le{}&
 A_\beta\Delta_{\mathrm E}(S,\Pi)+B_{\beta,\delta}b(S)
       -2b_{\mathrm{pin}}-2b_{\mathrm{both}}\\
 &+\frac{I(S,\Pi)-J(S)}{\ln2}.
 \end{aligned}
\]
The discovery construction is unchanged.
\end{lemma}
\begin{proof}
Each witness query contributes at most two hull caps, each of mass
less than $\delta w(S)$, so its newly supported mass is at most
$2\delta w(S)$. An appended terminal singleton lies in a positive
cap and obeys the same bound. Their masses sum to $b(S)$.

Apply Lemma~\ref{lem:capped-rank-credit} with
\[
 a_t=\varepsilon_\beta(t),\quad
 \tau=\varepsilon_\beta(N_\delta+1),\quad
 E=\{1,\ldots,N_\delta\},\quad \mu=0,\quad
 \zeta=1/\ln\beta.
\]
The discounts are nondecreasing. Its definition of $\chi$ is exactly
$\chi_\beta(\delta)$, since $s_\delta=1-2\delta N_\delta$.
Condition~\eqref{eq:flat-discrete-condition} is precisely
$\chi+\zeta\ge\tau$. Thus the positive-rank discounts and the
zero-route discount retained in
\eqref{eq:discrete-discovery-credit} save at least
$\chi_\beta(\delta)b(S)$.

The endpoint cost from Lemma~\ref{lem:pruning-macro} is unchanged.
Substituting this saving in its stage ledger replaces
$B_\beta$ by $B_{\beta,\delta}=B_\beta-\chi_\beta(\delta)$ and gives
the displayed bound. No witness masses are computed by the algorithm.
\end{proof}

The cap on each nonzero discovery class gives a useful entropy
lower bound, but the total mass of those classes matters as well.
They cannot all have mass $\delta$ unless their masses sum to the
positive mass $b(S)$; any remainder and the possible zero-pattern
singleton affect the least possible entropy. At total capped mass
one, the lowest-entropy arrangement fills as many entries to
$\delta$ as possible and leaves at most one smaller remainder;
$\sigma(\delta)$ is the entropy of that arrangement. The next lemma
shows that a smaller total capped mass still pays at least the same
entropy per unit mass, even when the remaining mass is placed in one
unrestricted entry.

\begin{lemma}[Entropy of capped discovery exits]
\label{lem:capped-discovery-entropy}
For $0<\delta<1$, put
\[
 k_\delta=\lfloor1/\delta\rfloor,\qquad
 r_\delta=1-k_\delta\delta,\qquad
 \sigma(\delta)=-k_\delta\delta\ln\delta-r_\delta\ln r_\delta,
\]
with $0\ln0=0$.
If a probability vector has one unrestricted entry $p_0$ and
all other entries at most $\delta$, its entropy is at least
$(1-p_0)\sigma(\delta)$.
\end{lemma}
\begin{proof}
Write $f(x)=-x\ln x$ with $f(0)=0$, and fix the capped total
$b=1-p_0$. Concavity of $f$ lets us transfer mass between any two
entries in $(0,\delta)$ until one is zero or the other is $\delta$,
without increasing entropy. Repeating gives the minimum-entropy shape:
$j=\lfloor b/\delta\rfloor$ entries of mass $\delta$ and at most
one remainder. Including the unrestricted entry, its entropy is
\[
 T(b)=f(1-b)+j f(\delta)+f(b-j\delta).
\]
It remains to show $T(b)\ge b\sigma(\delta)$.

On each interval between consecutive multiples of $\delta$, $T$ is
concave. It therefore suffices to prove the inequality at those
multiples and at the endpoints $0,1$. We have $T(0)=0$ and
$T(1)=\sigma(\delta)$. At a positive multiple $b=j\delta$,
\[
 \frac{T(b)}b=L+g(1-b),\qquad
 L=\ln(1/\delta),\qquad
 g(p)=\frac{-p\ln p}{1-p}.
\]
Since
$g'(p)=(-\ln p-1+p)/(1-p)^2\ge0$, these ratios decrease as the
multiple $b$ increases. Put $r=r_\delta$. If $r=0$, the last ratio
is $L=\sigma(\delta)$. If $r>0$, its difference from $T(1)$ is
\[
 L+\frac{f(r)}{1-r}-[(1-r)L+f(r)]
 =rL+\frac{r f(r)}{1-r}\ge0.
\]
Thus every required endpoint lies above the line
$b\sigma(\delta)$, and concavity supplies the intervals between them.
The given vector has entropy at least $T(b)$, proving the claim.
Also $\sigma(\delta)=\ln(1/\delta)+r_\delta
\ln(\delta/r_\delta)\ge\ln(1/\delta)$, with the second term
interpreted as zero when $r_\delta=0$.
\end{proof}

In a mass-flat state, every hull cap has mass less than
$\delta w(S)$. Each nonzero discovery class lies in its first
positive hull cap, or is an inferred singleton with a nonzero
pattern. Thus its mass is less than $\delta w(S)$.
Normalize the discovery-class masses by $w(S)$ and apply
Lemma~\ref{lem:capped-discovery-entropy},
using the zero-pattern class as the unrestricted entry. Multiplying
the resulting entropy bound by $w(S)$ gives
$J(S)\ge b(S)\sigma(\delta)$.
In particular, the zero-pattern mass is included rather than
discarded, and the capped distribution also accounts for any
remainder when the positive masses cannot all equal $\delta$.
Therefore
\begin{equation}\label{eq:pattern-entropy-slack}
 J(S)\ge b(S)\sigma(\delta)\ge b(S)\ln(1/\delta).
\end{equation}
If $\Sigma$ is the full pattern and $Z$ is the discovery transcript,
then $\Pi$ refines $\mathcal R$ and $Z$ is a function of $\Sigma$.
The chain rule of Lemma~\ref{lem:entropy-chain-rule} gives
\[
 I(S,\Pi)-J(S)=w(S)H(\Sigma\mid Z)\ge0,
\]
since conditional entropy is a weighted sum of terms
$-p\ln p\ge0$ for probabilities $p\in[0,1]$.
So $I(S,\Pi)\ge J(S)$.
The zero pattern, when present, is a singleton. Thus every
nonterminal class in a flat stage is a proper subset of $S$.

Let $T_{\mathrm E}(S)$ denote the tree returned by this recursion
when it is started on state $S$.

\begin{lemma}[Common stage-to-tree conversion]\label{lem:stage-conversion}
Suppose a recursive algorithm uses mass-balanced queries and flat
full-pattern stages, and every flat stage satisfies
\[
 W\le A\Delta_{\mathrm E}+Bb+\frac{I-J}{\ln2},\qquad
 0\le J\le I,\qquad J\ge b\sigma(\delta),
\]
with nonnegative $A,B$. Put
$c=\max\{1/h(\delta),B/\sigma(\delta)\}$. Then
\[
 c_{\mathrm E}(T(S))\le A\OPT_{\mathrm E}(S)+c\mathcal E(S)
 \le(A+c\ln2)\OPT_{\mathrm E}(S).
\]
\end{lemma}
\begin{proof}
Binary entropy is minimized on $[\delta,1-\delta]$ at an endpoint,
so a balanced query has information at least $w(S)h(\delta)$ and
costs at most $cI$. Also $c\ge1/\ln2$. A flat stage satisfies
\[
 Bb+\frac{I-J}{\ln2}\le cJ+c(I-J)=cI.
\]
Thus every stage costs at most $A\Delta_{\mathrm E}+cI$. The identities
$\Delta_{\mathrm E}=\OPT_{\mathrm E}(S)-\sum_A\OPT_{\mathrm E}(A)$
and $I=\mathcal E(S)-\sum_A\mathcal E(A)$ telescope over the disjoint
recursive classes; both potentials vanish at singleton leaves.
Lemma~\ref{lem:restriction} gives the last inequality.
\end{proof}

\begin{theorem}[Constant expected-depth approximation]
\label{thm:mass-average}
Let
\[
 h(\delta)=-\delta\ln\delta-(1-\delta)\ln(1-\delta),\qquad
 c(\delta)=\max\left\{\frac1{h(\delta)},
                    \frac{B_{\beta,\delta}}{\sigma(\delta)}\right\}.
\]
The chosen parameters satisfy \eqref{eq:flat-discrete-condition}.
For binary supplied-line queries and arbitrary positive hypothesis
probabilities, the preceding algorithm satisfies
\[
 c_{\mathrm E}(T_{\mathrm E}(S))\le
 A_\beta \OPT_{\mathrm E}(S)+c(\delta)\mathcal E(S)
 \quad\text{for every state }S.
\]
Its approximation factor is at most
\begin{equation}\label{eq:parameter-ratio}
 R_\beta(\delta)=A_\beta+
 \max\left\{\frac{\ln2}{h(\delta)},
            \frac{B_{\beta,\delta}\ln2}{\sigma(\delta)}\right\}.
\end{equation}
For the chosen rational threshold, this gives
\begin{equation}\label{eq:mass-average}
 \boxed{c_{\mathrm E}(T_{\mathrm E}(\HH))
         \le\rho \OPT_{\mathrm E}(\HH),\qquad
 \rho=A_\beta+\frac{B_{\beta,\delta}\ln2}{\sigma(\delta)}<5.445944.}
\end{equation}
Here $c_{\mathrm E}(T_{\mathrm E}(S))$ uses the original,
unnormalized probabilities, just like $\OPT_{\mathrm E}(S)$.
The approximation factor is independent of $n$ and of the smallest
hypothesis probability.
\end{theorem}
\begin{proof}
Lemma~\ref{lem:flat-discrete-credit}, with its nonpositive pinning terms
discarded, supplies the local stage bound. Equation~\eqref{eq:pattern-entropy-slack}
supplies the capped-transcript bound. Apply
Lemma~\ref{lem:stage-conversion} with $A=A_\beta$ and
$B=B_{\beta,\delta}$ to obtain the additive guarantee and
\eqref{eq:parameter-ratio}.

The concrete constant can be certified without floating-point
rounding. For rational $x\ge1$, put $z=(x-1)/(x+1)$ and
\[
 L_N(x)=2\sum_{j=0}^{N-1}\frac{z^{2j+1}}{2j+1}.
\]
The positive logarithm series and a geometric tail bound give
\begin{equation}\label{eq:rational-log-bounds}
 L_N(x)\le\ln x\le L_N(x)+
 \frac{2z^{2N+1}}{(2N+1)(1-z^2)}.
\end{equation}
Apply these rational bounds with $N=180$ to
$x=2$, $\beta$, $1/\delta$, $1/(1-\delta)$ and
$1/(64r_\delta)$. For the factorials below, range-reduce each
integer $j$ from $2$ through $35$ by writing
$j=2^k(j/2^k)$ with $j/2^k\in[1,2)$.
Use $\ln x=k\ln2+\ln(x/2^k)$ for this range reduction and
$\ln(1/r_\delta)=6\ln2+\ln(1/(64r_\delta))$.
The identity
$\ln(945/8)=3\ln3+\ln5+\ln7-3\ln2$ also avoids a slowly
converging logarithm series at the larger argument.
For integer $t$, set $s_t=\lfloor\beta t\rfloor$; the required
values of the discrete curve can be evaluated as
\[
 \Phi_\beta(t)=\frac{s_t+1}{2}
  +\frac{(t+s_t)\ln t-\ln(t!)-\ln(s_t!)}{2\ln\beta}.
\]
Compute $\varepsilon_\beta(2),\ldots,\varepsilon_\beta(5)$
from $t\le10$, expanding each log-factorial as a sum of logarithms.
Here $k_\delta=8$ and $N_\delta=4$. Substituting
$h(\delta)=\delta\ln(1/\delta)+(1-\delta)\ln(1/(1-\delta))$
gives the outward-rounded intervals
\[
 \begin{aligned}
 1.79590508&<\lambda<1.79590509,\qquad
 1.79590498<\nu<1.79590499,\\
 3.59181017&<A_\beta<3.59181018,\\
 0.736855718583&<\varepsilon_\beta(5)<0.736855718585,\\
 0.380956768686&<\chi_\beta(\delta)<0.380956768688,\\
 5.669579638912&<B_{\beta,\delta}<5.669579638914,\\
 2.119509588573&<\sigma(\delta)<2.119509588575,\\
 1.854133220728&<\frac{\ln2}{h(\delta)}<1.854133220730,\\
 1.854133221598&<\frac{B_{\beta,\delta}\ln2}{\sigma(\delta)}
               <1.854133221600.
 \end{aligned}
\]
This verifies $\lambda>\nu$,
\eqref{eq:flat-discrete-condition}, and dominance of the flat-stage
coefficient. In particular,
$\rho<3.59181018+1.854133221600<5.445944$.
\end{proof}

The two terms in \eqref{eq:parameter-ratio} describe the threshold
tradeoff. For fixed $\beta$, the balanced-query term decreases
and the flat-stage term increases. The mass constraint defining
$\chi_\beta(\delta)$ weakens as $\delta$ grows, so
$B_{\beta,\delta}$ is nondecreasing.
The capped entropy
$\sigma(\delta)$ is continuous and nonincreasing; between reciprocal
breakpoints its derivative is
$k_\delta\ln(r_\delta/\delta)\le0$.
The crossing of the two terms is the
minimizing threshold. The rational logarithm bounds certify
\[
 0.123498956866<\delta_*<0.123498956870.
\]
The chosen value gives $\rho\approx5.4459434004$ and lies within
$10^{-7}$ of that fixed-$\beta$ optimum. The parameters $\beta$
and $\delta$ were chosen together: $\beta$ is just above the
solution of $\beta-1=2\ln\beta$, where the positive-target and
zero-target coefficients satisfy $\lambda=\nu$.
The certified guarantee does not rely on floating-point comparisons.

\subsection{Implementation and running time}

Each stage computes a planar linear program, a convex hull, its
explicit pattern table, a minimum circular-interval cover and
the budgeted-coverage dynamic programs of
Lemma~\ref{lem:budget-coverage}. There are at most $|V|$ distinct
budgets and polynomially many shift intervals.
The two-support representation and the free opposite-edge cut
are computable by planar convex-hull operations.
Endpoint codes use polynomial interval dynamic programs.
Discovery's positive exits and terminal singleton are disjoint
subsets of $S$, so there are at most $|S|$ nonempty transcripts.
Each conditioned endpoint problem has at most $|S|$ possible
values. Thus the number and size of the dynamic programs are polynomial.

Each nonterminal stage partitions its state into at least two
nonempty proper children, so the recursive stage hierarchy has at
most $n-1$ internal stages. Constant questions, cached questions,
and empty branches are suppressed. The actual complete binary
decision tree then has $n-1$ internal nodes.
All arithmetic and comparisons have polynomial bit complexity.
In particular, the algorithm compares rational masses and optimizes
rational dynamic-program values; logarithms appear only in the
analysis. The fixed balance threshold is exactly
$1234989569/10^{10}$, and $\beta=3512863/10^6$ is rational as well.

The proof uses no lower bound on hypothesis probabilities.
A rare pattern may require many queries, but its small mass also
reduces its contribution to both accounting sums. This is why the
argument controls expected depth rather than maximum depth.
